\documentclass[10pt,twocolumn]{article}
\usepackage[margin=1in]{geometry}
\usepackage{booktabs}
\usepackage{hyperref}
\usepackage{graphicx}
\usepackage{amsmath}
\usepackage{amssymb}
\usepackage{microtype}
\usepackage{xcolor}
\usepackage{listings}
\usepackage{caption}
\usepackage{subcaption}

\hypersetup{
  colorlinks=true, linkcolor=blue, citecolor=blue, urlcolor=blue
}

\title{\textbf{Calibra: Dataset Observability and Coreset Selection\\
for Robot Imitation Learning}\\[0.5em]
\large Technical Report}

\author{
  Ömer Topçu\\
  \texttt{github.com/omerTT/Calibra}
}

\date{June 2026}

\begin{document}
\maketitle

\begin{abstract}
Robot learning labs routinely train policies on thousands of demonstration
episodes without inspecting the data. We present \textbf{Calibra}, an
open-source dataset observability toolkit that applies four deterministic
mathematical estimators — jerk spike rate, velocity discontinuity rate,
timing jitter, and log dimensionless jerk (LDLJ) — to detect quality
anomalies before training begins. We profile 15 public imitation-learning
datasets and find three actionable findings: (1) scripted motion planners
produce 22–25\% jerk spike rates, \emph{25$\times$} higher than human
teleoperation counterparts; (2) diverse hardware fleets (DROID) exhibit
10$\times$ more velocity discontinuities than single-platform datasets
(ALOHA); and (3) control frequency is a latent confounder that inflates
velocity discontinuity by up to 80 percentage points in coarse-grained
datasets (BridgeData~V2, 5~Hz). Calibra is pip-installable and runs on
any dataset with a LeRobot~v2 adapter in under two minutes.
\end{abstract}

% -----------------------------------------------------------------------
\section{Introduction}
\label{sec:intro}

Imitation learning pipelines treat data collection as a solved problem.
Demonstrations are recorded, concatenated into a dataset, and fed to a
training script. Quality inspection — if it happens at all — is manual,
inconsistent, and expensive.

In practice, real demonstration datasets contain three classes of silent
problems. \textbf{Mechanical noise:} jerk spikes from communication lag,
stiff joints, or abrupt operator corrections. \textbf{Temporal artifacts:}
dropped frames and duplicate timestamps from simulation engines under load.
\textbf{Distributional skew:} 60–80\% near-duplicate episodes in large
collections, inflating compute cost and reinforcing dominant modes.

None of these are detectable by standard loss curves or policy rollouts
during training. By the time a trained policy stalls or flails, the
practitioner has no way to determine whether the fault lies in the
architecture, the training recipe, or the data.

Calibra addresses the data side. It computes \emph{per-episode} and
\emph{per-dataset} diagnostic statistics, compares them against a
library of validated reference profiles, and flags anomalies with
bootstrap confidence intervals. Its four estimators are closed-form,
require no learned model, and run in $O(N)$ time over the episode length.

\paragraph{Contributions.}
\begin{enumerate}
  \item Four reproducible data-quality metrics with falsifiable claim
        thresholds validated across 15 public datasets.
  \item An empirical cross-dataset study revealing control-frequency
        confounding and hardware-diversity effects.
  \item An open-source toolkit (MIT licence) with a LeRobot~v2 adapter,
        a \texttt{compare} command, and a coreset pruner.
\end{enumerate}

% -----------------------------------------------------------------------
\section{Method}
\label{sec:method}

\subsection{Jerk Spike Rate}

For an episode of $T$ steps with action sequence
$\mathbf{a}_1,\ldots,\mathbf{a}_T$, the per-step jerk magnitude is
\[
  j_t = \left\| \mathbf{a}_{t+1} - 2\mathbf{a}_t + \mathbf{a}_{t-1} \right\|_2.
\]
A step is classified as a \emph{spike} if $j_t > 5 \cdot \mathrm{median}(j)$.
The episode-level spike rate is the fraction of steps classified as spikes.
This threshold is calibrated so that smooth 50~Hz ALOHA teleoperation produces
rates below 1\% and abrupt planner waypoints exceed 20\%.

\subsection{Velocity Discontinuity Rate}

\[
  \mathrm{vd}_t = \left\| \mathbf{a}_{t+1} - \mathbf{a}_t \right\|_2
    \big/ \max_s \left\| \mathbf{a}_s \right\|_2
\]
Steps where $\mathrm{vd}_t > 0.20$ count as discontinuities. The 20\%
threshold was chosen to distinguish genuine control-loop instabilities from
normal deceleration.

\subsection{Timing Jitter (CV)}

Given inter-step timestamps $\Delta t_1, \ldots, \Delta t_{T-1}$, the jitter
coefficient of variation is $\sigma_{\Delta t} / \mu_{\Delta t}$.  Values
above $10^{-2}$ indicate clock synchronisation problems; values above 0.1
indicate severe frame drops.

\subsection{Log Dimensionless Jerk (LDLJ)}

LDLJ~\cite{balasubramanian2012} is a scale-invariant smoothness index for
movement primitives:
\[
  \mathrm{LDLJ} = \log\!\left(
    \frac{D^3}{v_{\max}^2 \, T^5}
    \int_0^T \|\dddot{\mathbf{a}}\|^2 \, \mathrm{d}t
  \right)
\]
where $D$ is movement amplitude and $v_{\max}$ is peak speed. More-negative
values indicate less smooth motion. LDLJ is computed numerically from the
discrete action sequence via finite differences.

% -----------------------------------------------------------------------
\section{Experimental Setup}
\label{sec:experiments}

We profile 15 publicly available datasets, all accessible via the LeRobot~v2
API. Datasets span simulation and real hardware, human teleoperation and
scripted planners, and control frequencies from 5~Hz (BridgeData~V2) to
50~Hz (ALOHA). Table~\ref{tab:datasets} summarises the collection.

\begin{table}[htbp]
\centering
\small
\caption{Datasets profiled in this study.}
\label{tab:datasets}
\begin{tabular}{llrr}
\toprule
Dataset family & Mode & $N$ & Hz \\
\midrule
ALOHA sim human (×2)    & pos & 50 & 50 \\
ALOHA sim scripted (×2) & pos & 50 & 50 \\
ALOHA static (×4)       & pos & 49–50 & 50 \\
ALOHA mobile (×2)       & pos & 18–85 & 50 \\
DROID-100               & pos & 100 & 15 \\
BridgeData V2           & vel & 50\,415 & 5 \\
PushT / PushT-image     & vel & 206 & 10 \\
SVLA SO-100 (×2)        & pos & 50–56 & 50 \\
\bottomrule
\end{tabular}
\end{table}

All metric computations are deterministic and single-threaded; a 100-episode
dataset at 50~Hz takes under 15~seconds on a 2022 MacBook Pro.

% -----------------------------------------------------------------------
\section{Results}
\label{sec:results}

\subsection{Finding 1: Scripted planners create artifact noise (25$\times$)}

ALOHA simulation offers a controlled comparison: the \emph{same} bimanual
task executed by both human operators and a scripted planner at the same
frequency and action dimensionality. Human-collected episodes show a mean
jerk spike rate of 0.7--0.9\%; scripted episodes show 21.9--24.9\%. The
25$\times$ gap cannot be attributed to task difficulty — it reflects the
piecewise-linear trajectory profiles produced by the planner's
waypoint interpolation. Policies trained on scripted data may learn to
reproduce planner artifacts that do not appear in real deployment.

\subsection{Finding 2: Hardware diversity amplifies noise (10$\times$)}

DROID-100 collects 100 episodes across diverse hardware platforms, operator
backgrounds, and tasks. It shares the same control mode (position command)
as ALOHA hardware, but its mean jerk spike rate is 4.55\% and velocity
discontinuity is 7.08\% — roughly 10$\times$ the ALOHA hardware baselines
(0.42\% and 0.84\%). These are not pathological datasets; they represent
the realistic operating range of multi-robot, multi-operator data collection.
Calibra's \texttt{compare} command surfaces this gap automatically and
flags the 13 DROID episodes whose spike rate exceeds 8\% as outliers.

\subsection{Finding 3: Control frequency is a latent confounder}

BridgeData~V2 contains 50\,415 real-hardware demonstrations at 5~Hz.
Its velocity discontinuity rate is 80.5\% — far above any other
position- or velocity-command dataset we profiled. This is not hardware
noise: at 5~Hz, any smooth continuous trajectory appears as a sequence
of abrupt step changes when discretised, because the time resolution is
too coarse to capture within-step variation. This finding formally
falsified claim VD-002 ("velocity-command datasets: 10--20\%") and
prompted the frequency-qualified successor VD-003.

The lesson is methodological: \emph{velocity discontinuity rate is not
comparable across datasets without controlling for control frequency.}
Calibra's reference profiles record the control frequency alongside each
metric, and the \texttt{compare} command warns when the target dataset's
frequency differs by more than a factor of two from the reference.

\subsection{Finding 4: Timing quality is uniformly excellent}

Jitter CV is below $10^{-4}$ across all 15 datasets, and timestamp
dropout is zero in every profile. This suggests that LeRobot~v2's
normalised Parquet format provides reliable temporal alignment, and that
clock synchronisation issues documented in older Isaac Lab workflows are
absent in the published HuggingFace datasets.

\subsection{Quantitative summary}

Table~\ref{tab:cross_dataset} reports mean spike rate, velocity discontinuity,
jitter~CV, and LDLJ for all 15 datasets.

\begin{table*}[htbp]
\centering
\caption{Calibra observability metrics across 15 public imitation-learning datasets.
Scripted datasets are marked $\dagger$.
Spike and VelDisc are mean per-episode rates (\%).
LDLJ is mean log dimensionless jerk (more negative = less smooth).}
\label{tab:cross_dataset}
\small
\begin{tabular}{lllrrrr}
\toprule
Dataset & Ctrl & $N$ & Spike\,\% & VelDisc\,\% & Jitter CV & LDLJ \\
\midrule
aloha\_mobile\_cabinet          & position & 85     & 0.42  & 0.84  & 3.0e-05 & -24.70 \\
aloha\_mobile\_shrimp           & position & 18     & 0.46  & 0.66  & 8.5e-05 & -27.80 \\
aloha\_sim\_insertion\_human    & position & 50     & 0.69  & 2.39  & 1.1e-05 & -20.43 \\
aloha\_sim\_insertion\_scripted$^\dagger$ & position & 50 & 24.90 & 0.75 & 4.0e-06 & -17.19 \\
aloha\_sim\_transfer\_cube\_human & position & 50   & 0.94  & 4.56  & 4.0e-06 & -20.02 \\
aloha\_sim\_transfer\_cube\_scripted$^\dagger$ & position & 50 & 21.87 & 0.75 & 4.0e-06 & -17.23 \\
aloha\_static\_battery          & position & 49     & 0.44  & 5.46  & 1.4e-05 & -22.10 \\
aloha\_static\_candy            & position & 50     & 0.42  & 2.66  & 1.6e-05 & -22.28 \\
aloha\_static\_coffee           & position & 50     & 0.17  & 1.95  & 2.6e-05 & -24.14 \\
aloha\_static\_cups\_open       & position & 50     & 0.30  & 6.44  & 4.0e-06 & -20.51 \\
\midrule
BridgeData\_V2                  & velocity & 50{,}415 & 0.71 & \textbf{80.48} & 1.0e-06 & -13.79 \\
droid\_100                      & position & 100    & 4.55  & 7.08  & 6.0e-06 & -19.39 \\
pusht\_image                    & velocity & 206    & 7.76  & 11.11 & 3.0e-06 & -16.01 \\
pusht                           & velocity & 206    & 4.94  & 16.72 & 3.0e-06 & -16.34 \\
svla\_so100\_pickplace          & position & 50     & 1.09  & 6.73  & 9.0e-06 & -20.11 \\
svla\_so100\_stacking           & position & 56     & 1.02  & 1.65  & 9.0e-06 & -19.61 \\
\bottomrule
\end{tabular}
\end{table*}


% -----------------------------------------------------------------------
\section{Claim Registry}
\label{sec:claims}

Calibra maintains a falsifiable claim registry in
\texttt{calibra/claims/}. Each claim is a JSON record specifying the
metric, the dataset class, the assertion, a formal falsification
condition, and an evidence list. The registry currently contains 12
active claims across four metric families, backed by 15 reference
profiles (ratio $\geq 1$, per project policy).

Two claims have been formally falsified by the experiments in this
paper (VD-002, JS-002), demonstrating that the registry is live
empirical science rather than documentation. The falsification events
produced two successor claims (VD-003, JS-003) with refined scope.

% -----------------------------------------------------------------------
\section{Related Work}
\label{sec:related}

Dataset quality for robot learning has received growing attention.
Walke et al.~\cite{walke2023bridgedata} document BridgeData~V2 but do not report
control-smoothness metrics. Khazatsky et al.~\cite{khazatsky2024droid} profile DROID
for diversity but not for temporal or kinematic quality. Calibra provides
the first systematic cross-dataset comparison on the metrics that directly
affect policy learning stability.

Data-curation tools such as GRAPE~\cite{hejna2023} and
SOAR~\cite{ma2023} focus on reward-based episode selection; Calibra
is complementary, operating entirely on raw kinematic data without
requiring a reward model or environment access.

% -----------------------------------------------------------------------
\section{Conclusion}
\label{sec:conclusion}

We have shown that three systematic quality problems — planner artifacts,
hardware-diversity noise, and control-frequency confounding — are
detectable from raw kinematic data in under two minutes, before any
policy training takes place. Calibra's deterministic estimators and
falsifiable claim registry provide an auditable, reproducible quality
baseline for any dataset with LeRobot~v2 support.

Future work includes: (1) extending the adapter library to Isaac Lab and
ROS~2 bag formats; (2) integrating the coreset pruner with the
observability metrics to produce quality-weighted coresets; (3) running
a controlled ablation study measuring downstream policy performance as a
function of Calibra's quality scores.

% -----------------------------------------------------------------------
\section*{Acknowledgements}
This report was prepared with AI writing assistance (GitHub Copilot).
All experimental results, metric computations, and dataset profiles
are the sole work of the author.

\bibliographystyle{plain}
\begin{thebibliography}{9}

\bibitem{balasubramanian2012}
S.~Balasubramanian, A.~Melendez-Calderon, and E.~Burdet.
\newblock A robust and sensitive metric for quantifying movement smoothness.
\newblock \emph{IEEE Trans.\ Biomed.\ Eng.}, 59(8):2126--2136, 2012.

\bibitem{walke2023bridgedata}
H.~Walke et al.
\newblock BridgeData~V2: A dataset for robot learning at scale.
\newblock In \emph{CoRL}, 2023.

\bibitem{khazatsky2024droid}
A.~Khazatsky et al.
\newblock DROID: A large-scale in-the-wild robot manipulation dataset.
\newblock \emph{arXiv:2403.12945}, 2024.

\bibitem{hejna2023}
J.~Hejna and D.~Sadigh.
\newblock Few-shot preference learning for human-in-the-loop RL.
\newblock In \emph{CoRL}, 2023.

\bibitem{ma2023}
Y.~J.~Ma et al.
\newblock Eureka: Human-level reward design via coding large language models.
\newblock \emph{arXiv:2310.12931}, 2023.

\end{thebibliography}

\end{document}
